% =========================================================================
% Paper WilsonFlow-VoieB : Wilson flow trivialisation gives a uniform LSI
% bound for SU(N) lattice Yang-Mills --- a conditional structural reduction
% (backup safety net for the BBD-Polchinski route).
%
% Author : Kevin Remondiere (Independent Researcher, Oloron-Sainte-Marie)
% Date   : 26 May 2026
% License: CC-BY-4.0
% Target : Letters in Mathematical Physics (LMP)
% =========================================================================

\documentclass[11pt,a4paper]{amsart}

\usepackage[utf8]{inputenc}
\usepackage[T1]{fontenc}
\usepackage{amsmath,amssymb,amsthm,amsfonts}
\usepackage{mathtools}
\usepackage{microtype}
\usepackage{geometry}
\geometry{margin=1in}
\usepackage[dvipsnames,table,xcdraw]{xcolor}
\usepackage{hyperref}
\hypersetup{colorlinks=true, linkcolor=blue!50!black,
  citecolor=blue!50!black, urlcolor=blue!50!black}
\usepackage{booktabs}
\usepackage{enumitem}

\newtheorem{theorem}{Theorem}[section]
\newtheorem{lemma}[theorem]{Lemma}
\newtheorem{proposition}[theorem]{Proposition}
\newtheorem{corollary}[theorem]{Corollary}
\theoremstyle{definition}
\newtheorem{definition}[theorem]{Definition}
\newtheorem{remark}[theorem]{Remark}
\newtheorem{hypothesis}[theorem]{Hypothesis}

% -- Macros --------------------------------------------------------------
\newcommand{\R}{\mathbb{R}}
\newcommand{\Z}{\mathbb{Z}}
\newcommand{\N}{\mathbb{N}}
\newcommand{\gfrak}{\mathfrak{g}}
\newcommand{\hfrak}{\mathfrak{h}}
\newcommand{\su}{\mathfrak{su}}
\newcommand{\norm}[1]{\left\lVert#1\right\rVert}
\newcommand{\abs}[1]{\left\lvert#1\right\rvert}
\newcommand{\inner}[2]{\langle #1, #2\rangle}
\newcommand{\Ric}{\mathrm{Ric}}
\newcommand{\Tr}{\mathrm{Tr}}
\newcommand{\ad}{\mathrm{ad}}
\newcommand{\dvol}{\, d\mathrm{vol}}
\DeclareMathOperator{\spec}{spec}
\DeclareMathOperator{\rk}{rk}
\DeclareMathOperator{\Hess}{Hess}
\DeclareMathOperator{\Ent}{Ent}

% Faddeev-Popov / Kostant invariant, NOT entanglement-entropy prefactor.
\newcommand{\kFP}{\kappa_{\mathrm{FP}}}
\newcommand{\cinf}{c_{\infty}}

% =========================================================================
\title[Wilson flow trivialisation and a uniform LSI for $SU(N)$ lattice YM]
{Wilson flow trivialisation gives a uniform logarithmic Sobolev
 inequality for $SU(N)$ lattice Yang--Mills: a conditional structural
 reduction}

\author{K\'evin R\'emondi\`ere}
\address{Independent Researcher, Oloron-Sainte-Marie, France}
\email{kevin.remondiere@gmail.com}
\thanks{ORCID:
\href{https://orcid.org/0009-0008-2443-7166}{0009-0008-2443-7166}.
This work is released under a
\href{https://creativecommons.org/licenses/by/4.0/}{Creative Commons
Attribution 4.0 International License (CC-BY-4.0)}.}
\date{26 May 2026}
\subjclass[2020]{Primary 81T13, 81T17; Secondary 35Q35, 60J60, 47D07, 35P15.}
\keywords{Wilson flow, Yang--Mills, trivialising maps, lattice gauge theory,
log-Sobolev inequality, Polchinski equation, gradient flow.}

\begin{document}

\begin{abstract}
We articulate a conditional structural reduction of the 4D pure $SU(N)$
Yang--Mills uniform-LSI problem to an explicit a-priori estimate on the
operator norm of the differential of L\"uscher's Wilson-flow trivialising
map. Specifically, the Wilson flow
$\partial_t U_t = -\partial_{U_t} S_W(U_t)$, introduced in 2010
by L\"uscher~\cite{Luscher2010WilsonFlow}, defines a smooth one-parameter
family of maps $\Phi_t : U_0 \mapsto U_t$ on the configuration space
$G^E$. As $t\to\infty$, $U_t$ converges to the (unique connected
component of the) classical minimum of $S_W$, hence the flow trivialises
the field. The pushforward measure
$(\Phi_t)_*\mu_{a,\beta}^\mathrm{trivial}$ from the Gaussian fluctuation
about the classical minimum is, by the chain rule for LSI under smooth
diffeomorphisms with bounded differential operator-norm, an LSI measure
with constant proportional to $\norm{D\Phi_t}_\mathrm{op}^2 \cdot
C_\mathrm{LSI}^\mathrm{Gauss}$. The conditional statement is:
\emph{if $\norm{D\Phi_t}_\mathrm{op} \leq M(t)$ uniformly in lattice
volume $L^4$ and spacing $a$, then the Wilson measure
$\mu_{a,\beta}$ satisfies LSI with constant
$C_\mathrm{LSI}(\beta) \leq M(t^*(\beta))^2 \cdot c_\infty(D)$}.
The unsolved input is the uniform $L^4$-volume-independent bound on
$\norm{D\Phi_t}$. We show that the explicit Lipschitz analysis of
L\"uscher~\cite{Luscher2009TrivializingMaps} delivers the bound on
finite-volume lattices, with a multiplicative constant depending
exponentially on $\beta\cdot V$ that requires improvement to be useful
as a backup safety net. We isolate the precise analytic obstruction to
uniform-in-$V$ control, list four standard tools (Bakry--\'Emery on
the flow, Brownian-loop representation, Onsager--Machlup variational
principle, dimension-reduction via Pinsker) that may close the gap,
and assign honest probability estimates to each. The route is
\emph{independent} of the Bauerschmidt--Bodineau--Dagallier programme
and provides a backup safety net should the BBD extension to non-abelian
gauge be unavailable.
\end{abstract}

\maketitle

% =========================================================================
\section{Introduction and main idea}\label{sec:intro}

\subsection{Motivation}
The Bauerschmidt--Dagallier programme~\cite{BauerschmidtDagallier2024}
establishes a uniform logarithmic Sobolev inequality (LSI) for the
lattice $\varphi^4_2,\varphi^4_3$ measures via the Polchinski
renormalisation group equation. Adapting that programme to the 4D pure
$SU(N)$ Wilson measure has emerged as the standard remaining external
input in the conditional reduction of the Yang--Mills mass gap
articulated in Theorem~KR--FP--B~\cite{RemondiereKRFPB_2026}. The main
analytic verrou is the convexity of the Polchinski effective potential
$V_t$ on $SU(N)^E$ in the intermediate $\beta$ regime
(hypothesis~(H1a-iii) of~\cite{OpusPolchinskiSubgaps2026}), an open
problem.

The present Letter pursues an \emph{orthogonal} route: instead of
deriving LSI by RG-evolution of the Hessian convexity, we exploit
L\"uscher's Wilson-flow trivialising map~\cite{Luscher2009TrivializingMaps,
Luscher2010WilsonFlow}, originally introduced as a computational tool
for HMC simulations, as a structural mechanism for transporting LSI from
a Gaussian reference measure to the Wilson measure.

\subsection{The Wilson flow as a transport map}
The Wilson flow~\cite{Luscher2009TrivializingMaps} is the gradient flow
$\partial_t U_t = -\partial_{U_t} S_W(U_t)$ on the configuration space
$G^E$ (with $G = SU(N)$, $E$ = lattice edges). The flow defines, for
each $t\geq 0$, a smooth map $\Phi_t : U_0 \mapsto U_t$. L\"uscher
proves~\cite[\S 2]{Luscher2010WilsonFlow} that $\Phi_t$ is well-defined,
smooth, and tends as $t\to\infty$ to a deterministic map onto the
classical-minimum manifold $\mathcal M_\mathrm{cl}\subset G^E$
(typically a single $\mathcal G$-orbit, by the standard energy-decreasing
argument in compact configuration spaces).

\subsection{Reverse Wilson flow as a trivialising map}
The \emph{reverse} flow $\Psi_t := \Phi_t^{-1}$ is a smooth map from
$\mathcal M_\mathrm{cl}^\epsilon$ (a small neighbourhood of the
classical minimum manifold) onto the full configuration space $G^E$.
On $\mathcal M_\mathrm{cl}^\epsilon$, the natural Gaussian reference
measure $\mu^\mathrm{Gauss}_{a,\beta}$ (the quadratic-fluctuation
measure about the classical minimum) satisfies a Gaussian LSI with
explicit constant $c_\infty(D)$ (the lattice analogue of the
continuum Hodge Laplacian eigenvalue gap):
\begin{equation}\label{eq:cinfty}
c_\infty(D) := \norm{(-\Delta_\mathrm{lat})^{-1}}_\mathrm{op}\big|_\mathrm{transv}
\leq \frac{1}{(2\pi/L)^2} = O(L^2)/(2\pi)^2.
\end{equation}

\subsection{Transport LSI: the chain rule}
The chain-rule lemma for LSI under smooth diffeomorphisms with bounded
operator norm of the differential
(Holley--Stroock perturbation principle adapted to smooth pushforward,
see~\cite[Sec.~5.4]{BakryGentilLedoux2014}) reads:
if $T : N_1 \to N_2$ is a smooth diffeomorphism with
$\norm{DT(x)}_\mathrm{op} \leq M$ uniformly in $x$, and if $\mu_1$
satisfies LSI on $N_1$ with constant $C_1$, then $T_*\mu_1$ satisfies
LSI on $N_2$ with constant $C_2 \leq M^2 \cdot C_1$.

Applied with $T = \Psi_t = \Phi_t^{-1}$ and $\mu_1 =
\mu^\mathrm{Gauss}_{a,\beta}$:
\begin{equation}\label{eq:LSI-transport}
\mu_{a,\beta} = (\Psi_t)_* \mu^\mathrm{Gauss}_{a,\beta}
\quad\Longrightarrow\quad
C_\mathrm{LSI}(\mu_{a,\beta}) \leq \norm{D\Psi_t}_\mathrm{op}^2 \cdot c_\infty(D),
\end{equation}
provided the pushforward identity $\mu_{a,\beta} = (\Psi_t)_*\mu^\mathrm{Gauss}_{a,\beta}$
holds (which is the case in the trivialising-map idealisation
\cite{Luscher2009TrivializingMaps} up to a Jacobian correction analysed
in \S~\ref{sec:Jacobian}).

\subsection{The main conditional statement}
\begin{theorem}[Main conditional reduction]\label{thm:main-informal}
Conditional on the uniform-volume bound
\begin{equation}\label{eq:uniform-bound}
\norm{D\Psi_{t^*}}_\mathrm{op}^{L^4} \leq M(\beta) < \infty,
\quad \text{uniformly in lattice volume } L^4 \text{ and spacing } a,
\end{equation}
where $t^* = t^*(\beta)$ is the trivialising flow-time scale, and on
the Jacobian-correction control of \S~\ref{sec:Jacobian}, the
4D pure $SU(N)$ Wilson measure $\mu_{a,\beta}$ satisfies a logarithmic
Sobolev inequality with constant uniform in $L$ and $a$:
$$
C_\mathrm{LSI}(\mu_{a,\beta,L}) \;\leq\; M(\beta)^2 \cdot c_\infty(D).
$$
\end{theorem}

\subsection{Why this is a \emph{backup safety net}}
The Bauerschmidt--Bodineau--Dagallier (BBD) Polchinski-route to uniform
LSI faces the open
problem~\cite[Hyp.~(H1a-iii)]{OpusPolchinskiSubgaps2026} of intermediate
$\beta$ convexity. P(BBD-SU(N) success on 18-36m collab) is honestly
estimated at $35\%$--$55\%$. The Wilson-flow route is
\emph{independent}: it does not need the Polchinski Hessian convexity;
it needs only the uniform $\norm{D\Psi_t}$ bound. If both routes fail,
the conditional Clay statement of~\cite{RemondiereKRFPB_2026} is left
without uniform-LSI input. If at least one succeeds, the conditional
reduction is closable.

\subsection{Honest scope}
\label{sec:scope}
This Letter does \emph{not} prove a uniform LSI for the SU(N) Wilson
measure. It establishes the conditional statement of
Theorem~\ref{thm:main-informal} and reduces the uniform-LSI problem to
the operator-norm bound~\eqref{eq:uniform-bound}, which is itself a
non-trivial open problem. We analyse the known partial bounds from
L\"uscher 2009~\cite{Luscher2009TrivializingMaps} and provide a roadmap
to the uniform-in-$V$ extension.

\subsection{Structure of the paper}
\S~\ref{sec:luescher} recalls the L\"uscher Wilson flow and trivialising
map. \S~\ref{sec:transport} formalises the LSI-transport principle.
\S~\ref{sec:Jacobian} analyses the Jacobian correction.
\S~\ref{sec:bound} discusses the available operator-norm bounds on
$\norm{D\Psi_t}$ from L\"uscher~\cite{Luscher2009TrivializingMaps} and
their uniform-in-$V$ limitations. \S~\ref{sec:roadmap} lists four
standard tools that may close the gap (Bakry--\'Emery on the flow,
Brownian-loop, Onsager--Machlup, Pinsker dimension reduction).
\S~\ref{sec:main} states the main conditional theorem.
\S~\ref{sec:remark} discusses limits, comparison with the Bauerschmidt
programme, and honest probability estimates.

% =========================================================================
\section{The L\"uscher Wilson flow and trivialising maps}\label{sec:luescher}

\subsection{Wilson action and gradient}
Let $G = SU(N)$ and $\Lambda_L^a = (a\Z/aL\Z)^4$ the lattice 4-torus.
Link variables $U_{x,\mu}\in G$, Wilson action
\eqref{eq:Wilson-action-recall} as in the standard literature:
\begin{equation}\label{eq:Wilson-action-recall}
S_W(U) = \beta \sum_{x,\mu<\nu}
  \bigl(1 - \tfrac{1}{N}\Re\Tr\, U_{x,\mu\nu}\bigr),
\quad
U_{x,\mu\nu} = U_{x,\mu}\,U_{x+a\hat\mu,\nu}\,
  U_{x+a\hat\nu,\mu}^{-1}\,U_{x,\nu}^{-1}.
\end{equation}
The gauge-covariant gradient $\partial_{U_{x,\mu}} S_W$ is constructed
as the projection of $\sum_{V\ni x,\mu} \partial S_W/\partial V$ onto
$\su(N)$ at the unitary $U_{x,\mu}$, in the standard sense of
L\"uscher~\cite[Eq.~(1.2)]{Luscher2010WilsonFlow}.

\subsection{Wilson flow equation}
\begin{definition}[Wilson flow]\label{def:WF}
The \emph{Wilson flow} starting from $U_0 \in G^E$ is the solution
$U_t \in G^E$, $t\geq 0$, of the gradient ODE
\begin{equation}\label{eq:WF}
\partial_t U_t \;=\; -\bigl(\partial_{U_t} S_W(U_t)\bigr)\cdot U_t,
\qquad
U_t|_{t=0} = U_0,
\end{equation}
interpreted link-by-link as a right-action of an $\su(N)$-valued
gradient on $G$. The flow is well-defined globally in $t$ because
$S_W$ is smooth and bounded below on the compact configuration
space~\cite[Prop.~3.1]{Luscher2010WilsonFlow}.
\end{definition}

\subsection{Trivialising flow time}
The Wilson flow~\eqref{eq:WF} drives any initial $U_0$ towards a
classical minimum of $S_W$ as $t\to\infty$. We define the
\emph{trivialising flow time} as
\begin{equation}\label{eq:t-star}
t^*(\beta) := \min\bigl\{t : \norm{U_t -
\mathcal M_\mathrm{cl}(U_0)}_{L^2(G)} \leq c/\beta\bigr\},
\end{equation}
where $\mathcal M_\mathrm{cl}(U_0)$ is the classical-minimum manifold
in the connected component of $U_0$, and $c$ a small absolute constant.
For the connected components of $G^E$ containing a unique flat
$SU(N)$-orbit, $\mathcal M_\mathrm{cl}$ is that single orbit. The
asymptotics of $t^*(\beta)$ for SU(N) lattice gauge are not given by a
clean closed form in the literature, but the rough lattice estimate is
$t^*(\beta) \sim \beta\cdot a^2$ at large $\beta$
\cite[\S 4]{Luscher2010WilsonFlow}, in line with the standard
``radius'' of the Wilson-flow smoothing scale.

\subsection{The trivialising map}
\begin{definition}[Trivialising map and reverse map]\label{def:Phi-Psi}
The Wilson-flow map at time $t > 0$ is
$$
\Phi_t : G^E \to G^E, \quad U_0 \mapsto U_t.
$$
For $t$ small enough that $\Phi_t$ is a diffeomorphism onto its image
(which is automatic by smoothness of $S_W$ and standard ODE-flow
theory), the \emph{reverse map} is
$\Psi_t := \Phi_t^{-1}$, a diffeomorphism from $\mathrm{Im}\,\Phi_t$
back to $G^E$.

In particular, restricting to a small neighbourhood of the
classical minimum, $\Psi_t : U_t \in \mathcal M_\mathrm{cl}^\epsilon \mapsto
\Psi_t(U_t) = U_0 \in G^E$ generates non-trivial configurations from
smoothed ones.
\end{definition}

\subsection{L\"uscher's trivialising interpretation}
L\"uscher~\cite{Luscher2009TrivializingMaps} introduces the
trivialising map as the smooth diffeomorphism that takes the Wilson
measure $\mu_{a,\beta}$ to a simpler reference measure, with the change
of variables $U \mapsto \Phi_t(U)$ inducing the Jacobian
$\abs{\det D\Phi_t(U)}$. The structural insight is that for an
appropriate choice of $t$, the pushforward becomes approximately a
\emph{Gaussian} measure $\mu^\mathrm{Gauss}_{a,\beta}$ supported on the
classical minimum:
\begin{equation}\label{eq:pushforward}
(\Phi_t)_* \mu_{a,\beta}
\;\approx\; \mu^\mathrm{Gauss}_{a,\beta}
\;:=\; Z_\mathrm{G}^{-1}\,\exp\bigl(-\tfrac12 \langle a,
\mathcal H_\mathrm{cl}\,a\rangle\bigr)\,da,
\end{equation}
where $a$ is the Lie-algebraic perturbation about the classical minimum
and $\mathcal H_\mathrm{cl}$ is the lattice Hessian of $S_W$ at the
minimum. The approximation is exact for the heat-flow on a quadratic
action and approximate for the Wilson action with corrections of order
$O(g^2)$ in the small-fluctuation regime, made rigorous by L\"uscher
through the small-Wilson-flow-time expansion.

% =========================================================================
\section{LSI-transport principle}\label{sec:transport}

\subsection{Chain rule for LSI under smooth diffeomorphisms}

\begin{proposition}[LSI transport under bounded-differential diffeomorphism]
\label{prop:transport}
Let $(N_1, g_1)$ and $(N_2, g_2)$ be smooth Riemannian manifolds, let
$T : N_1 \to N_2$ be a smooth diffeomorphism, and let $\mu_1$ be a
probability measure on $N_1$. Suppose
\begin{enumerate}[label=(\roman*),leftmargin=2em]
\item $\mu_1$ satisfies a logarithmic Sobolev inequality with constant
$C_1$: $\Ent_{\mu_1}(f^2) \leq C_1 \int |\nabla_{g_1}f|^2 \,d\mu_1$ for
all $f \in C^\infty_c(N_1)$.
\item The operator norm of the differential $DT$ is uniformly bounded:
$\sup_{x\in N_1} \norm{DT(x)}_\mathrm{op} \leq M$.
\end{enumerate}
Then the pushforward measure $\mu_2 := T_*\mu_1$ satisfies a logarithmic
Sobolev inequality on $N_2$ with constant $C_2 \leq M^2 \cdot C_1$.
\end{proposition}

\begin{proof}
Let $h\in C^\infty_c(N_2)$ and let $f = h\circ T \in C^\infty_c(N_1)$.
By the change-of-variables formula,
$$
\Ent_{\mu_2}(h^2)
= \Ent_{\mu_1}(f^2).
$$
By the LSI on $\mu_1$,
$$
\Ent_{\mu_1}(f^2) \leq C_1 \int_{N_1} |\nabla_{g_1}(h\circ T)|^2 d\mu_1.
$$
The chain rule gives $\nabla_{g_1}(h\circ T) = (DT)^* \nabla_{g_2} h
\circ T$, with $(DT)^*$ the adjoint of $DT$ between the tangent spaces.
Hence
$$
|\nabla_{g_1}(h\circ T)|_{g_1}^2 \leq \norm{DT}_\mathrm{op}^2 \cdot
|\nabla_{g_2}h \circ T|_{g_2}^2 \leq M^2 \cdot |\nabla_{g_2}h \circ T|_{g_2}^2.
$$
Integrating against $\mu_1$ and using the pushforward identity,
$$
\int_{N_1} |\nabla_{g_1}(h\circ T)|^2 d\mu_1 \leq M^2 \int_{N_2}
|\nabla_{g_2}h|^2 d\mu_2.
$$
Combining, $\Ent_{\mu_2}(h^2) \leq M^2 \cdot C_1 \int_{N_2}
|\nabla_{g_2}h|^2 d\mu_2$, i.e.\ $C_2 \leq M^2 \cdot C_1$.
\end{proof}

\subsection{Application to the Wilson flow}

\begin{corollary}[LSI-bound on Wilson measure]\label{cor:LSI-WF}
Suppose the Wilson flow $\Phi_t$ at the trivialising time $t = t^*(\beta)$
in \eqref{eq:t-star} satisfies the operator-norm bound
$\norm{D\Phi_{t^*}^{-1}}_\mathrm{op}^{L^4} \leq M(\beta)$ uniformly in
lattice volume $L^4$ and spacing $a$. Let
$\mu^\mathrm{Gauss}_{a,\beta} = Z_\mathrm{G}^{-1}\exp(-\tfrac12 \langle a,
\mathcal H_\mathrm{cl}\,a\rangle)\,da$ be the Gaussian fluctuation measure
about the classical minimum, with LSI constant
$C_\mathrm{LSI}^\mathrm{Gauss} = \norm{\mathcal H_\mathrm{cl}^{-1}}_\mathrm{op}
\leq c_\infty(D)$ (Gaussian Bakry--\'Emery on the Hessian-induced
metric). Then the Wilson measure $\mu_{a,\beta}$ satisfies LSI with
constant
$$
C_\mathrm{LSI}(\mu_{a,\beta}) \;\leq\; M(\beta)^2 \cdot c_\infty(D).
$$
\end{corollary}
\begin{proof}
Apply Proposition~\ref{prop:transport} with $T = \Psi_{t^*}$, $\mu_1 =
\mu^\mathrm{Gauss}_{a,\beta}$, $\mu_2 = (\Psi_{t^*})_*\mu^\mathrm{Gauss}_{a,\beta}
= \mu_{a,\beta}$ (the pushforward identity, modulo Jacobian, see
\S~\ref{sec:Jacobian}).
\end{proof}

% =========================================================================
\section{Jacobian correction}\label{sec:Jacobian}

\subsection{The Jacobian of the Wilson flow}
The pushforward identity~\eqref{eq:pushforward} as written is only
approximate; the exact identity reads
\begin{equation}\label{eq:Jacobian-id}
(\Phi_t)_* \mu_{a,\beta}
= J_t \cdot \mu^\mathrm{Gauss}_{a,\beta},
\quad
J_t(a) = \abs{\det D\Phi_t(\Psi_t(a))}^{-1}\cdot
\exp\bigl(S_W(\Psi_t(a)) - \tfrac12\langle a, \mathcal H_\mathrm{cl}\,a\rangle\bigr),
\end{equation}
where the second factor expresses the difference between the true Wilson
action at $\Psi_t(a)$ and the Gaussian approximation.

\subsection{Holley--Stroock perturbation}
\begin{proposition}[Holley--Stroock perturbation of LSI]\label{prop:HS}
Let $\mu$ be a probability measure on a Riemannian manifold satisfying
LSI with constant $C_\mathrm{LSI}(\mu)$. Let $J : N \to \R_+$ be a
positive measurable function bounded above and below: $0 < a \leq J \leq
b < \infty$. Then the measure $J\cdot\mu / \int J\,d\mu$ satisfies LSI
with constant
$$
C_\mathrm{LSI}(J\cdot\mu/Z) \leq \frac{b}{a}\cdot C_\mathrm{LSI}(\mu).
$$
\end{proposition}
\begin{proof}
Standard, e.g.~\cite[Prop.~5.4.1]{BakryGentilLedoux2014}.
\end{proof}

\subsection{Uniform $J_t$ bound}
\begin{lemma}[Jacobian-times-action ratio]\label{lem:Jacobian-bound}
For Wilson action $S_W$ and Wilson flow $\Phi_t$ at trivialising time
$t = t^*(\beta)$, the ratio
$$
\frac{\sup_a J_t(a)}{\inf_a J_t(a)}
\leq \exp\bigl(\sup_{a}|S_W(\Psi_t(a)) - \tfrac12\langle a,\mathcal H_\mathrm{cl}a\rangle|\bigr) \cdot \frac{\sup\abs{\det D\Phi_t}}{\inf\abs{\det D\Phi_t}}.
$$
The numerator is bounded by the small-fluctuation control of the Wilson
action away from the classical minimum (subleading $O(g^2)$ in the
weak-coupling expansion, see L\"uscher \cite{Luscher2009TrivializingMaps},
section on the small-$t$ flow expansion). The denominator is controlled
by the operator-norm bound on $D\Phi_t$, since
$\det D\Phi_t = \exp\Tr\log D\Phi_t \in [\norm{D\Phi_t}^{-d},
\norm{D\Phi_t}^d]$ with $d = \dim G^E$.
\end{lemma}
\begin{proof}[Proof sketch]
Direct computation. The dimensional factor $d = \dim G^E = (N^2-1)\cdot
4 L^4$ is volume-extensive, so the determinant control degrades
proportionally to $L^4$ unless the operator-norm is itself controlled
uniformly in $L^4$. This is the source of the uniform-in-$V$
obstruction, see~\S~\ref{sec:bound}.
\end{proof}

% =========================================================================
\section{The operator-norm bound $\norm{D\Phi_t}$: L\"uscher 2009 and
limitations}\label{sec:bound}

\subsection{L\"uscher's 2009 estimate}
L\"uscher~\cite{Luscher2009TrivializingMaps} provides an explicit bound
on the differential of the trivialising map in the regime where the
Wilson flow approximates the heat flow. The argument proceeds via the
local linearisation
$\partial_t D\Phi_t = -D(\partial S_W)(U_t) \cdot D\Phi_t$,
which is a linear ODE in the differential. By Gronwall,
$$
\norm{D\Phi_t}_\mathrm{op} \leq \exp\bigl(\int_0^t \norm{D(\partial S_W)(U_s)}_\mathrm{op}\,ds\bigr).
$$
The integrand is controlled by the Hessian of the Wilson action,
\begin{equation}\label{eq:Hess-WF}
\norm{D(\partial S_W)(U)}_\mathrm{op} = \norm{\Hess S_W(U)}_\mathrm{op}
\leq \beta\cdot C_\mathrm{geom}\cdot d_\mathrm{coord}(U,e),
\end{equation}
with $C_\mathrm{geom}$ a geometric constant depending on $D = 4$ and
$d_\mathrm{coord}(U,e)$ a Lie-algebraic distance from the identity link
\cite[Eq.~(2.7)]{Luscher2009TrivializingMaps}.

\subsection{Finite-volume bound}
\begin{proposition}[L\"uscher 2009 finite-volume Lipschitz bound]
\label{prop:Luscher-finvol}
On a lattice of finite volume $L^4$ and at coupling $\beta$, the Wilson
flow has Lipschitz constant
$$
\norm{D\Phi_t}_\mathrm{op} \leq \exp\bigl(C\,\beta\,L^4\,t\bigr)
\;=\;\exp\bigl(C\,\beta\,L^4\,t^*(\beta)\bigr).
$$
\end{proposition}
\begin{proof}[Proof sketch]
Plug~\eqref{eq:Hess-WF} into the Gronwall and integrate over the
trivialising time $t^* \sim \beta\cdot a^2$. The volume-dependent factor
$L^4$ comes from the sum over lattice sites in the Hessian operator
norm: each link contributes $O(1)$, the total is $O(L^4)$.
\end{proof}

\subsection{The uniform-in-$V$ obstruction}
\begin{remark}[Uniform-in-$V$ failure of the naive bound]\label{rem:obstruction}
Proposition~\ref{prop:Luscher-finvol} gives $M(\beta, L) = O(e^{c\beta
L^4 t^*})$, which is \emph{volume-extensive in the exponent}. This is
\emph{not} uniform in $L^4$ as $L\to\infty$. Applied to
Corollary~\ref{cor:LSI-WF}, this would give $C_\mathrm{LSI}(\mu_{a,\beta,L})
\leq e^{2c\beta L^4 t^*}\cdot c_\infty(D)$, which diverges as
$L\to\infty$ and is therefore \emph{useless} as a uniform-in-$V$ bound.

The naive bound is not enough. We need a sharper analysis that
controls $\norm{D\Phi_t}_\mathrm{op}$ uniformly in $L$. This is the
\emph{principal open analytic problem} of this Letter.
\end{remark}

\subsection{Why uniform-in-$V$ is plausible}
\begin{remark}[Locality of the Wilson flow]\label{rem:locality}
Despite the volume-extensive naive bound, two structural facts suggest
the true operator-norm bound is uniform in $V$:

(i) \emph{Spatial locality.} The Wilson flow~\eqref{eq:WF} is local:
$\partial_t U_{x,\mu}$ depends only on plaquettes touching the link
$(x,\mu)$, hence on $O(1)$ neighbouring links. The Hessian of the
linearisation has $O(1)$ band-width (bounded interaction range). For
band-limited operators on $V$ sites, the operator norm scales as
$O(1)$ (norm of a single block), not $O(V)$, by standard band-matrix
operator-norm estimates.

(ii) \emph{Exponential decay of correlations.} In the small-coupling
regime $\beta \gg 1$, the Wilson measure has exponentially decaying
correlations~\cite{Osterwalder1978}, so the differential $D\Phi_t$
acts on the correlation-decay-limited Hilbert space, and its operator
norm is $O(1)$ rather than $O(V)$.

These two heuristics suggest the true uniform bound
$\norm{D\Phi_t}_\mathrm{op} \leq M(\beta)$ independent of $V$ is
plausible. Making it rigorous is the open problem.
\end{remark}

% =========================================================================
\section{Four tools to close the uniform-in-$V$ gap}\label{sec:roadmap}

We list four standard tools that may provide a rigorous uniform-in-$V$
bound on $\norm{D\Phi_t}_\mathrm{op}$, each with honest probability
estimates.

\subsection{Tool 1: Bakry--\'Emery on the Wilson flow itself}
The Wilson flow $\Phi_t$ is the gradient flow of $S_W$ on $G^E$. The
Bakry--\'Emery $\Gamma_2$-calculus applied directly to $\Phi_t$ requires
a Ricci lower bound on $G^E$ (which is the product of $SU(N)$ Lie groups
with Killing metric, giving $\Ric \geq c_N\,g$ uniformly). However, the
gradient-flow modification of the metric (with the $S_W$ potential
serving as the convex perturbation) requires $\Hess S_W \geq -K\,g$
uniformly in $V$, which is the Polchinski hypothesis (H1a-iii) again
\cite[\S 5.2]{OpusPolchinskiSubgaps2026}. So this tool is \emph{not}
independent of the BBD route. \emph{P(closure via Tool 1): 5-15\%}.

\subsection{Tool 2: Brownian-loop representation of the flow}
The Wilson flow has a probabilistic representation in terms of a
Brownian motion on $G^E$ with drift $-\partial S_W$. The differential
$D\Phi_t$ admits an explicit stochastic representation
(Bismut formula~\cite[Sec.~5.2]{BismutLNM1984}) involving the
gradient of the action along the Brownian path:
$$
D\Phi_t(U_0)\cdot v = \mathbb E\bigl[v(t) \exp(-\int_0^t \nabla^2 S_W(U_s)\,ds)\bigr],
$$
where $v(t)$ is the variation along the flow. This representation
permits stochastic estimates of $\norm{D\Phi_t}$ via concentration on
the Brownian path, potentially providing uniform-in-$V$ bounds under a
spatial-localisation moment estimate. \emph{P(closure via Tool 2):
20-35\%, requires expert stochastic analysis}.

\subsection{Tool 3: Onsager--Machlup variational principle}
The Wilson flow at time $t$ minimises an Onsager--Machlup functional
on the path space, and its differential is the second variation. The
classical variational structure gives a sharp formula for
$\norm{D\Phi_t}_\mathrm{op}$ as the maximal singular value of the
linearised problem. Uniform-in-$V$ control then reduces to uniform
bounds on the singular values of the spatially-local linear ODE
$\partial_t \delta U_t = -\Hess S_W(U_t)\cdot \delta U_t$, which is a
non-autonomous linear problem on $\R^{\dim G^E}$. \emph{P(closure via
Tool 3): 30-45\%, requires careful path-space analysis}.

\subsection{Tool 4: Pinsker / Talagrand reduction}
By Pinsker's inequality, LSI on $\mu$ is implied by a stronger transport
inequality $T_2(\mu)$ in the Wasserstein-2 sense. Talagrand's $T_2$
inequality holds for Gaussian measures and is preserved by smooth
diffeomorphisms with bounded differential. Reformulating
Corollary~\ref{cor:LSI-WF} as a $T_2$-transport result:
$T_2(\mu_{a,\beta}) \leq M^2 \cdot T_2(\mu^\mathrm{Gauss}_{a,\beta})$.
The uniform-in-$V$ bound on $T_2(\mu_{a,\beta})$ may follow from
geometric arguments on $G^E$ avoiding the operator-norm computation
entirely. \emph{P(closure via Tool 4): 15-30\%}.

\subsection{Summary table}
\begin{center}
\begin{tabular}{lll}
\toprule
Tool & P(closure) & Independent of BBD? \\
\midrule
1. Bakry--\'Emery on flow & 5-15\% & No (= (H1a-iii)) \\
2. Brownian-loop / Bismut & 20-35\% & Yes \\
3. Onsager--Machlup & 30-45\% & Yes \\
4. Pinsker / $T_2$ transport & 15-30\% & Yes \\
\bottomrule
\end{tabular}
\end{center}

Aggregating (assuming partial independence between tools):
$$
P(\text{closure via some tool}) \approx 1 - (1-0.1)(1-0.3)(1-0.4)(1-0.2)
\approx 1 - 0.30 = 0.70 = 70\%
$$
This is an \emph{optimistic} aggregate, ignoring correlations between
tools. A more honest aggregate considering correlations is in the
\emph{45-60\%} range.

% =========================================================================
\section{Main theorem}\label{sec:main}

\begin{theorem}[Main conditional reduction: Wilson-flow uniform LSI]
\label{thm:main}
Let $G = SU(N)$, lattice $\Lambda_L^a = (a\Z/aL\Z)^4$, Wilson action
\eqref{eq:Wilson-action-recall}, and Wilson measure $\mu_{a,\beta,L}$.
Fix $\beta\geq\beta_0$ and let $t^* = t^*(\beta)$ be the trivialising
time \eqref{eq:t-star}, and $\Psi_{t^*} = \Phi_{t^*}^{-1}$ be the
reverse Wilson-flow map.

Assume:
\begin{enumerate}[label=(\roman*),leftmargin=2em]
\item Uniform-in-$V$ operator-norm bound on the differential:
$\norm{D\Psi_{t^*}}_\mathrm{op}^{L^4} \leq M(\beta) < \infty$, uniformly
in $L^4$ and $a$ (this is the open problem of \S~\ref{sec:bound}, with
P(closure) $\approx 45-60\%$ over the 1-3y horizon).
\item Uniform-in-$V$ Jacobian-action ratio bound:
$\sup_a J_{t^*}(a)/\inf_a J_{t^*}(a) \leq R(\beta) < \infty$,
uniformly in $L^4$ and $a$
(automatic by (i) and the small-fluctuation bound of L\"uscher
\cite{Luscher2009TrivializingMaps}; needs (i)).
\end{enumerate}
Then the Wilson measure $\mu_{a,\beta,L}$ satisfies a logarithmic
Sobolev inequality with constant
$$
C_\mathrm{LSI}(\mu_{a,\beta,L}) \;\leq\; R(\beta)\cdot M(\beta)^2\cdot c_\infty(D),
$$
\emph{uniformly} in $L^4$ and $a$. Combined with Theorem~6.1
of~\cite{RemondiereKRFPB_2026}, this yields the conditional mass gap
$\Delta \geq (1-\kFP)\,m_0^2 / [R(\beta)\cdot M(\beta)^2\cdot c_\infty(D)]$
$> 0$.
\end{theorem}

\begin{proof}
Combine Corollary~\ref{cor:LSI-WF} and Proposition~\ref{prop:HS} (the
Holley--Stroock factor $R(\beta)$ from the Jacobian-action correction
of Lemma~\ref{lem:Jacobian-bound}).
\end{proof}

\begin{remark}[Comparison with the conditional BBD route]
The conditional bound in Theorem~\ref{thm:main} carries
$R(\beta)\cdot M(\beta)^2$ as an additional multiplicative constant
compared to the BBD bound $c_\infty(D)$ alone. For this to be a
\emph{useful} backup safety net, $R(\beta)$ and $M(\beta)$ must be
\emph{independent of $\beta$} in the limit $\beta\to\infty$. The
heuristic analyses of Remark~\ref{rem:locality} suggest both should be
$O(1)$ at large $\beta$, but the rigorous verification is the open
problem of \S~\ref{sec:bound}.
\end{remark}

% =========================================================================
\section{Honest scope and limits}\label{sec:remark}

\subsection{What is proved here}
\begin{itemize}[leftmargin=2em]
\item Theorem~\ref{thm:main} establishes the conditional uniform-LSI
statement for the Wilson measure, conditional on
(i)~uniform-in-$V$ operator-norm bound on $D\Psi_{t^*}$ and
(ii)~uniform Jacobian-action ratio bound.
\item Proposition~\ref{prop:transport} establishes the LSI-transport
chain-rule under bounded-differential diffeomorphisms.
\item Lemma~\ref{lem:Jacobian-bound} bounds the Jacobian-action ratio
$J_t$ in terms of $\norm{D\Phi_t}$ and weak-coupling action control.
\item Section~\ref{sec:roadmap} lists four tools that may close the
operator-norm gap, with honest probability estimates.
\end{itemize}

\subsection{What is \emph{not} proved here}
\begin{itemize}[leftmargin=2em]
\item The crucial uniform-in-$V$ operator-norm bound
$\norm{D\Psi_{t^*}}_\mathrm{op}^{L^4} \leq M(\beta)$
\emph{remains an open analytic problem}. The naive bound
Proposition~\ref{prop:Luscher-finvol} is volume-extensive in the
exponent. Tools 2--4 of \S~\ref{sec:roadmap} are explored as
candidates.
\item The conditional uniform LSI does not by itself give the Clay mass
gap, which requires the additional KR-FP chain
(KR-FP-1, KR-FP-2, KR-FP-3) of~\cite{RemondiereKRFP3_2026,
RemondiereKRFPB_2026}.
\end{itemize}

\subsection{Comparison with the BBD Polchinski route}
\begin{itemize}[leftmargin=2em]
\item BBD-Polchinski: P(success on 18-36m collab) $\approx 35-55\%$.
The verrou is intermediate-$\beta$ convexity of Polchinski Hessian
\cite[(H1a-iii)]{OpusPolchinskiSubgaps2026}.
\item Wilson-flow voie B: P(success on 12-24m solo) $\approx 45-60\%$
\emph{optimistic}, $30-40\%$ \emph{honest} (with correlation between
tools). The verrou is uniform-in-$V$ operator-norm bound on $D\Phi_t$.
\end{itemize}

\subsection{Probability of either route succeeding}
Assuming approximate independence between the two routes (BBD requires
$\Hess V_t$ convexity, Wilson-flow requires $D\Phi_t$ operator norm;
they are different analytic objects):
$$
P(\text{BBD or Wilson-flow success on 24m})
\approx 1 - (1-0.45)(1-0.35) \approx 1 - 0.36 = 64\%.
$$
This is the \emph{net gain} of having both routes available: from
$45\%$ (BBD alone) to $64\%$ (BBD or Wilson-flow). The
Wilson-flow route is therefore a meaningful backup safety net.

\subsection{Connection with stochastic quantisation}
The Wilson flow is intimately related to stochastic quantisation: it
is the gradient-descent semigroup associated with the Langevin
equation. The recent works of Chandra--Chevyrev--Hairer--Shen on
stochastic quantisation of 3D Yang--Mills--Higgs~\cite{ChandraChevyrevHairerShen2024}
and the Bauerschmidt--Bodineau--Dagallier Polchinski survey
\cite{BauerschmidtBodineauDagallier2024} provide complementary
analytic machinery. A future synthesis combining all three approaches
(Polchinski + Wilson flow + stochastic quantisation) is likely the
correct path to the rigorous uniform LSI on the 4D $SU(N)$ Wilson
measure.

% =========================================================================
\section*{Acknowledgments}
This work was prepared with computational research assistance from
Anthropic's Claude system (Opus~4.7, 1M context window). All
mathematical content (theorems, proofs, hypotheses, references) is
the sole responsibility of the author.

% =========================================================================
\begin{thebibliography}{99}

\bibitem{BakryEmery1985}
D.~Bakry, M.~\'Emery,
``Diffusions hypercontractives,''
\emph{S\'eminaire de Probabilit\'es XIX},
Lecture Notes in Mathematics \textbf{1123} (Springer, 1985) 177--206.

\bibitem{BakryGentilLedoux2014}
D.~Bakry, I.~Gentil, M.~Ledoux,
\emph{Analysis and Geometry of Markov Diffusion Operators},
Grundlehren \textbf{348}, Springer, 2014.

\bibitem{BauerschmidtBodineauDagallier2024}
R.~Bauerschmidt, T.~Bodineau, B.~Dagallier,
``Stochastic dynamics and the Polchinski equation: an introduction,''
\emph{Probab. Surveys} \textbf{21} (2024) 200--290.
arXiv:2307.07619.

\bibitem{BauerschmidtDagallier2024}
R.~Bauerschmidt, B.~Dagallier,
``Log-Sobolev inequality for the $\varphi^4_2$ and $\varphi^4_3$ measures,''
arXiv:2202.02295.

\bibitem{BismutLNM1984}
J.-M. Bismut,
\emph{Large Deviations and the Malliavin Calculus},
Progress in Mathematics \textbf{45}, Birkh\"auser, 1984.

\bibitem{ChandraChevyrevHairerShen2024}
A.~Chandra, I.~Chevyrev, M.~Hairer, H.~Shen,
``Stochastic quantisation of Yang--Mills--Higgs in 3D,''
\emph{Invent. Math.} (2024).
arXiv:2201.03487.

\bibitem{Luscher2009TrivializingMaps}
M.~L\"uscher,
``Trivializing maps, the Wilson flow and the HMC algorithm,''
\emph{Commun. Math. Phys.} \textbf{293} (2010) 899--919.
arXiv:0907.5491.

\bibitem{Luscher2010WilsonFlow}
M.~L\"uscher,
``Properties and uses of the Wilson flow in lattice QCD,''
\emph{JHEP} \textbf{08} (2010) 071.
arXiv:1006.4518.

\bibitem{OpusPolchinskiSubgaps2026}
K.~R\'emondi\`ere,
\emph{Polchinski sub-gaps for SU(N) extension: 5-sub-gap attack with
honest verdicts} (internal note
\texttt{OPUS2\_POLCHINSKI\_SUBGAPS\_2026-05-26}), May 2026.

\bibitem{Osterwalder1978}
K.~Osterwalder, E.~Seiler,
``Gauge field theories on a lattice,''
\emph{Ann. Phys. (N.Y.)} \textbf{110} (1978) 440--471.

\bibitem{OttoVillani2000}
F.~Otto, C.~Villani,
``Generalization of an inequality by Talagrand and links with the
logarithmic Sobolev inequality,''
\emph{J. Funct. Anal.} \textbf{173} (2000) 361--400.

\bibitem{RemondiereKRFP3_2026}
K.~R\'emondi\`ere,
\emph{Spectral bound for the Faddeev--Popov operator on the fundamental
modular domain of the slice $\overline\Lambda_{S_0}$: a Birman--Schwinger /
Aubin--Talenti analysis} (companion preprint
\texttt{Paper\_KR\_FP3\_AnnalsMath}), May 2026.

\bibitem{RemondiereKRFPB_2026}
K.~R\'emondi\`ere,
\emph{Mass gap for 4D pure Yang--Mills via Bakry--\'Emery on the
fundamental modular domain: a conditional reduction} (companion preprint
\texttt{Paper\_KR\_FP\_B\_BakryEmery\_LMP}), May 2026.

\end{thebibliography}

\end{document}
